%sección III
\section{Trabajos Relacionados}\label{sec:Trabajos-Relacionados}

La revisión de literatura se realizó a partir de artículos científicos publicados entre
2021 y 2025, consultados en las bases de datos ACM, EBSCO, IET y Springer Nature Link. Los
trabajos seleccionados abordan aspectos relacionados con la seguridad en la computación
en la nube, su aplicación en entornos educativos y el diseño de arquitecturas seguras.

Aljumah y Ahanger analizan los principales riesgos y mecanismos de
defensa en la computación en la nube, destacando la necesidad de implementar controles
robustos para garantizar la protección de la información. En el ámbito educativo,
Chatterjee et al. proponen un modelo de seguridad para sistemas de
gestión de aprendizaje en la nube, enfocado en la confidencialidad de los datos y la
aceptación de los usuarios.

Por otro lado, Al-Sharif y Bushnag presentan un enfoque basado en
aprendizaje automático para la detección de intrusiones en entornos de nube, evidenciando
el potencial de técnicas adaptativas para enfrentar amenazas avanzadas. Complementariamente,
Ntentos et al. abordan la seguridad desde la fase de diseño de
arquitecturas en la nube mediante Infraestructura como Código, resaltando la importancia
de prácticas seguras y automatizadas.

Finalmente, los trabajos de Yang y Yasser et al. analizan
el impacto general de la computación en la nube en contextos educativos y administrativos,
aunque con un enfoque principalmente conceptual y sin profundizar en la implementación
de controles de seguridad.

En conjunto, la literatura revisada evidencia la ausencia de estudios orientados a la
definición de una arquitectura de seguridad institucional alineada con las normas
ISO/IEC 27001, 27017 y 27018, particularmente en entornos universitarios de investigación,
lo cual fundamenta el desarrollo del presente trabajo.
